% ===== PRÉAMBULE CANONIQUE — à coller VERBATIM en tête de CHAQUE .tex =====
\documentclass[11pt,a4paper]{article}
\usepackage[utf8]{inputenc}\usepackage[T1]{fontenc}\usepackage{lmodern}
\usepackage[french]{babel}
\usepackage{amsmath,amssymb,mathtools}\usepackage{icomma}
\usepackage{siunitx}\sisetup{locale=FR}  % \qty{}{}, \si{}, \ang{}
\usepackage{xcolor}\usepackage{colortbl}
\usepackage{tikz}\usepackage{pgfplots}\pgfplotsset{compat=1.18}
\usepackage[siunitx,europeanresistors]{circuitikz} % circuits (élec.)
\usepackage[most]{tcolorbox}\usepackage{enumitem}\usepackage{array,booktabs}
\usepackage[a4paper,margin=2.2cm]{geometry}
\usetikzlibrary{arrows.meta,calc,angles,quotes,positioning,patterns,%
  backgrounds,shapes.geometric,decorations.pathmorphing,decorations.markings,intersections}
\definecolor{primary}{RGB}{222,49,99}\definecolor{primarydark}{RGB}{150,22,55}
\definecolor{defcol}{RGB}{0,114,178}\definecolor{methcol}{RGB}{52,92,148}
\definecolor{excol}{RGB}{0,135,90}\definecolor{attncol}{RGB}{200,40,55}
\tcbset{boxbase/.style={enhanced,breakable,boxrule=0.9pt,arc=2.5pt,
  left=3mm,right=3mm,top=2.3mm,bottom=2.3mm,fonttitle=\bfseries,coltitle=white}}
% --- boîtes TUTORIEL (noms EXACTS, ne pas renommer) ---
\newtcolorbox{cle}[1][]{boxbase,colback=defcol!6,colframe=defcol,
  title={\textsc{À retenir}\ifx\relax#1\relax\relax\else\textmd{ : \itshape #1}\fi}}
\newtcolorbox{methode}[1][]{boxbase,colback=methcol!7,colframe=methcol,sharp corners,
  title={\textsc{Méthode}\ifx\relax#1\relax\relax\else\textmd{ --- \itshape #1}\fi}}
\newtcolorbox{exemplebox}[1][]{boxbase,colback=excol!5,colframe=excol,
  title={\textsc{Exemple}\ifx\relax#1\relax\relax\else\textmd{ : \itshape #1}\fi}}
\newtcolorbox{piege}[1][]{boxbase,colback=attncol!6,colframe=attncol,
  title={\textsc{Piège}\ifx\relax#1\relax\relax\else\textmd{ : \itshape #1}\fi}}
\newtcolorbox{remarque}[1][]{boxbase,colback=black!4,colframe=black!45,
  title={\textsc{Remarque}\ifx\relax#1\relax\relax\else\textmd{ : \itshape #1}\fi}}
% --- boîtes EXERCICE / CORRIGÉ (noms EXACTS, auto-comptées) ---
\newtcolorbox[auto counter]{exo}[1][]{boxbase,colback=primary!4,colframe=primary,
  title={\textsc{Exercice}~\thetcbcounter\ifx\relax#1\relax\relax\else\textmd{ --- \itshape #1}\fi}}
\newtcolorbox[auto counter]{corrige}[1][]{boxbase,colback=excol!4,colframe=excol,
  title={\textsc{Corrigé}~\thetcbcounter\ifx\relax#1\relax\relax\else\textmd{ --- \itshape #1}\fi}}
\newcommand{\vv}[1]{\vec{#1}}\newcommand{\ds}{\displaystyle}
\newcommand{\R}{\mathbb{R}}\newcommand{\N}{\mathbb{N}}\newcommand{\Z}{\mathbb{Z}}
% (un \usepackage{fancyhdr}+en-tête est possible ; il est ignoré côté web)
% =========================================================================
\begin{document}
\sisetup{per-mode=symbol}

\begin{exo}[traversée d'une rivière]
Un bateau dirige son étrave perpendiculairement à la rive et avance à \qty{4}{\metre\per\second}.
La rivière, large de \qty{60}{\metre}, a un courant de \qty{3}{\metre\per\second} parallèle aux rives.
\begin{enumerate}[label=\alph*)]
  \item Quelle est la vitesse du bateau par rapport au sol ?
  \item De quel angle sa trajectoire réelle est-elle déviée ?
  \item Combien de temps met-il pour atteindre l'autre rive ?
  \item De quelle distance a-t-il été emporté vers l'aval ?
\end{enumerate}
\end{exo}

\begin{exo}[décomposition inverse]
Une balle part avec une vitesse de norme \qty{20}{\metre\per\second}. On souhaite que sa composante
verticale soit $v_y=\qty{10}{\metre\per\second}$.
\begin{enumerate}[label=\alph*)]
  \item Sous quel angle $\theta$ faut-il la lancer ?
  \item Que vaut alors sa composante horizontale $v_x$ ?
\end{enumerate}
\end{exo}

\begin{exo}[addition par composantes]
Deux vitesses perpendiculaires se composent : $\vv{v}_1=\qty{5}{\metre\per\second}$ vers l'Est et
$\vv{v}_2=\qty{12}{\metre\per\second}$ vers le Nord.
\begin{enumerate}[label=\alph*)]
  \item Quelle est la norme de la vitesse résultante ?
  \item Quel angle fait-elle avec la direction Est ?
\end{enumerate}
\end{exo}

\begin{exo}[accélérer à norme constante]
Un point décrit un cercle à vitesse de norme constante \qty{2}{\metre\per\second} (MCU). On compare le
vecteur vitesse en deux points séparés d'un quart de tour (les directions font alors un angle de \ang{90}).
\begin{enumerate}[label=\alph*)]
  \item Représente $\vv{v}_1$ et $\vv{v}_2$ et détermine la norme de la variation $\Delta\vv{v}=\vv{v}_2-\vv{v}_1$.
  \item Que peut-on conclure sur l'accélération, alors que la norme de la vitesse n'a pas changé ?
\end{enumerate}
\end{exo}

\begin{exo}[avion dans le vent]
Un avion vole vers le Nord à \qty{160}{\kilo\metre\per\hour} par rapport à l'air. Un vent d'Ouest
(soufflant vers l'Est) de \qty{120}{\kilo\metre\per\hour} le pousse perpendiculairement.
\begin{enumerate}[label=\alph*)]
  \item Quelle est la vitesse de l'avion par rapport au sol ?
  \item De quel angle sa trajectoire est-elle déviée par rapport au Nord ?
\end{enumerate}
\end{exo}

\end{document}
